
\Needspace{12\baselineskip}
\subsubsection{Operadores de las \(13\) multisecciones}

Las cifras de familias consignadas en cada ficha son cardinalidades locales del
soporte de esa multisección. No se suman entre fichas: el cociente de las
56 familias de ruta y el cociente de las 13 multisecciones son dos particiones
distintas del atlas de 81 celdas, relacionadas por los mapas de proyección y no
por una descomposición disjunta común.

\begingroup
\small
\setlength{\parskip}{.12\baselineskip}
\setlist[itemize]{leftmargin=1.65em,itemsep=0pt,topsep=.1em,
  parsep=0pt,partopsep=0pt}
\let\texttt\textnormal
Las fichas siguientes publican el operador en su codominio nativo. Cada espectro se reproduce íntegramente; ninguna lista de autovalores se sustituye por una coordenada de masa.
\Needspace{10\baselineskip}
\noindent\textbf{sector leptónico cargado, profundidad G0 (centro).} \(\mathcal B_{\ell,\mathrm{G0}}\).\par
\noindent\textbf{Dominio de fibra.} Rango 1; celda 41; sin órbita cromática; 1 familia de ruta.
\noindent\textbf{Espectro completo de \(K_D\).}
\begin{itemize}
\item \texttt{\detokenize{(80,54,6)x1}}
\end{itemize}
\noindent\textbf{Espectro completo de \(K_\Omega\).}
\begin{itemize}
\item \texttt{\detokenize{(80,54,6)x1}}
\end{itemize}
\noindent\textbf{Conclusión en la fibra.} Operador diagonal exacto sobre la fibra \(B_1\); la ruta nominal y el morfismo declarado publican el carácter escalar sin sustituir el espectro de la fibra.
\medskip
\Needspace{10\baselineskip}
\noindent\textbf{sector leptónico cargado, profundidad G2.} \(\mathcal B_{\ell,\mathrm{G2}}\).\par
\noindent\textbf{Dominio de fibra.} Rango 8; celdas 21,23,25,39,43,57,59,61; sin órbita cromática; 8 familias de ruta.
\noindent\textbf{Espectro completo de \(K_D\).}
\begin{itemize}
\item \texttt{\detokenize{(0,24,0)x1}}
\item \texttt{\detokenize{(40,78,27)x2}}
\item \texttt{\detokenize{(80,54,7)x1}}
\item \texttt{\detokenize{(80,66,2)x1}}
\item \texttt{\detokenize{(80,66,3)x1}}
\item \texttt{\detokenize{(100,66,0)x1}}
\item \texttt{\detokenize{(100,66,2)x1}}
\end{itemize}
\noindent\textbf{Espectro completo de \(K_\Omega\).}
\begin{itemize}
\item \texttt{\detokenize{(80,54,6)x8}}
\end{itemize}
\noindent\textbf{Conclusión en la fibra.} Operador diagonal exacto sobre la fibra \(B_1\); la ruta nominal y el morfismo declarado publican el carácter escalar sin sustituir el espectro de la fibra.
\medskip
\Needspace{22\baselineskip}
\noindent\textbf{sector leptónico cargado, profundidad G4.} \(\mathcal B_{\ell,\mathrm{G4}}\).\par
\noindent\textbf{Dominio de fibra.} Rango 16; celdas 1,3,5,7,9,19,27,37,45,55,63,73,75,77,79,81; sin órbita cromática; 16 familias de ruta.
\noindent\textbf{Espectro completo de \(K_D\).}
\begin{itemize}
\item \texttt{\detokenize{(0,24,0)x1}}
\item \texttt{\detokenize{(40,24,2)x2}}
\item \texttt{\detokenize{(40,54,6)x2}}
\item \texttt{\detokenize{(40,84,30)x2}}
\item \texttt{\detokenize{(60,78,25)x3}}
\item \texttt{\detokenize{(80,66,2)x2}}
\item \texttt{\detokenize{(80,66,3)x3}}
\item \texttt{\detokenize{(80,66,4)x1}}
\end{itemize}
\noindent\textbf{Espectro completo de \(K_\Omega\).}
\begin{itemize}
\item \texttt{\detokenize{(24,54,2)x1}}
\item \texttt{\detokenize{(56,54,5)x4}}
\item \texttt{\detokenize{(80,54,6)x11}}
\end{itemize}
\noindent\textbf{Conclusión en la fibra.} Operador diagonal exacto sobre la fibra \(B_1\); la ruta nominal y el morfismo declarado publican el carácter escalar sin sustituir el espectro de la fibra.
\medskip
\Needspace{10\baselineskip}
\noindent\textbf{sector neutrínico, profundidad G1.} \(\mathcal B_{\nu,\mathrm{G1}}\).\par
\noindent\textbf{Dominio de fibra.} Rango 2; celdas 40,42; sin órbita cromática; 2 familias de ruta.
\noindent\textbf{Espectro completo de \(K_D\).}
\begin{itemize}
\item \texttt{\detokenize{(40,24,2)x1}}
\item \texttt{\detokenize{(40,78,27)x1}}
\end{itemize}
\noindent\textbf{Espectro completo de \(K_\Omega\).}
\begin{itemize}
\item \texttt{\detokenize{(40,54,4)x1}}
\item \texttt{\detokenize{(80,54,6)x1}}
\end{itemize}
\noindent\textbf{Conclusión en la fibra.} Operador diagonal exacto sobre la fibra \(B_1\); la ruta nominal y el morfismo declarado publican el carácter escalar sin sustituir el espectro de la fibra.
\medskip
\Needspace{10\baselineskip}
\noindent\textbf{sector neutrínico, profundidad G2.} \(\mathcal B_{\nu,\mathrm{G2}}\).\par
\noindent\textbf{Dominio de fibra.} Rango 4; celdas 22,24,58,60; sin órbita cromática; 4 familias de ruta.
\noindent\textbf{Espectro completo de \(K_D\).}
\begin{itemize}
\item \texttt{\detokenize{(0,24,0)x1}}
\item \texttt{\detokenize{(40,84,30)x1}}
\item \texttt{\detokenize{(60,78,25)x1}}
\item \texttt{\detokenize{(80,66,3)x1}}
\end{itemize}
\noindent\textbf{Espectro completo de \(K_\Omega\).}
\begin{itemize}
\item \texttt{\detokenize{(4,54,2)x1}}
\item \texttt{\detokenize{(80,54,6)x3}}
\end{itemize}
\noindent\textbf{Conclusión en la fibra.} Operador diagonal exacto sobre la fibra \(B_1\); la ruta nominal y el morfismo declarado publican el carácter escalar sin sustituir el espectro de la fibra.
\medskip
\Needspace{18\baselineskip}
\noindent\textbf{sector neutrínico, profundidad G3.} \(\mathcal B_{\nu,\mathrm{G3}}\).\par
\noindent\textbf{Dominio de fibra.} Rango 6; celdas 20,26,38,44,56,62; sin órbita cromática; 6 familias de ruta.
\noindent\textbf{Espectro completo de \(K_D\).}
\begin{itemize}
\item \texttt{\detokenize{(40,24,2)x2}}
\item \texttt{\detokenize{(40,78,27)x1}}
\item \texttt{\detokenize{(60,84,28)x1}}
\item \texttt{\detokenize{(80,54,6)x1}}
\item \texttt{\detokenize{(80,66,3)x1}}
\end{itemize}
\noindent\textbf{Espectro completo de \(K_\Omega\).}
\begin{itemize}
\item \texttt{\detokenize{(36,54,4)x1}}
\item \texttt{\detokenize{(56,54,5)x1}}
\item \texttt{\detokenize{(80,54,6)x4}}
\end{itemize}
\noindent\textbf{Conclusión en la fibra.} Operador diagonal exacto sobre la fibra \(B_1\); la ruta nominal y el morfismo declarado publican el carácter escalar sin sustituir el espectro de la fibra.
\medskip
\Needspace{10\baselineskip}
\noindent\textbf{sector neutrínico, profundidad G4.} \(\mathcal B_{\nu,\mathrm{G4}}\).\par
\noindent\textbf{Dominio de fibra.} Rango 8; celdas 2,4,6,8,74,76,78,80; sin órbita cromática; 8 familias de ruta.
\noindent\textbf{Espectro completo de \(K_D\).}
\begin{itemize}
\item \texttt{\detokenize{(0,24,0)x1}}
\item \texttt{\detokenize{(40,24,2)x1}}
\item \texttt{\detokenize{(40,78,27)x2}}
\item \texttt{\detokenize{(40,84,30)x1}}
\item \texttt{\detokenize{(80,54,7)x1}}
\item \texttt{\detokenize{(80,66,2)x1}}
\item \texttt{\detokenize{(100,66,0)x1}}
\end{itemize}
\noindent\textbf{Espectro completo de \(K_\Omega\).}
\begin{itemize}
\item \texttt{\detokenize{(36,54,4)x1}}
\item \texttt{\detokenize{(56,54,5)x1}}
\item \texttt{\detokenize{(80,54,6)x6}}
\end{itemize}
\noindent\textbf{Conclusión en la fibra.} Operador diagonal exacto sobre la fibra \(B_1\); la ruta nominal y el morfismo declarado publican el carácter escalar sin sustituir el espectro de la fibra.
\medskip
\Needspace{10\baselineskip}
\noindent\textbf{rama quark de tipo inferior, profundidad G1.} \(\mathcal B_{d,\mathrm{G1}}\).\par
\noindent\textbf{Dominio de fibra.} Rango 2; celdas 32,50; B | G; 2 familias de ruta.
\noindent\textbf{Espectro completo de \(K_D\).}
\begin{itemize}
\item \texttt{\detokenize{(40,24,2)x1}}
\item \texttt{\detokenize{(40,78,27)x1}}
\end{itemize}
\noindent\textbf{Espectro completo de \(K_\Omega\).}
\begin{itemize}
\item \texttt{\detokenize{(40,54,4)x1}}
\item \texttt{\detokenize{(80,54,6)x1}}
\end{itemize}
\noindent\textbf{Conclusión en la fibra.} Operador diagonal exacto sobre la fibra \(B_1\); la ruta nominal y el morfismo declarado publican el carácter escalar sin sustituir el espectro de la fibra.
\medskip
\Needspace{10\baselineskip}
\noindent\textbf{rama quark de tipo inferior, profundidad G2.} \(\mathcal B_{d,\mathrm{G2}}\).\par
\noindent\textbf{Dominio de fibra.} Rango 4; celdas 30,34,48,52; B | G | R; 4 familias de ruta.
\noindent\textbf{Espectro completo de \(K_D\).}
\begin{itemize}
\item \texttt{\detokenize{(0,24,0)x1}}
\item \texttt{\detokenize{(40,84,30)x1}}
\item \texttt{\detokenize{(60,78,25)x1}}
\item \texttt{\detokenize{(80,66,2)x1}}
\end{itemize}
\noindent\textbf{Espectro completo de \(K_\Omega\).}
\begin{itemize}
\item \texttt{\detokenize{(4,54,2)x1}}
\item \texttt{\detokenize{(80,54,6)x3}}
\end{itemize}
\noindent\textbf{Conclusión en la fibra.} Operador diagonal exacto sobre la fibra \(B_1\); la ruta nominal y el morfismo declarado publican el carácter escalar sin sustituir el espectro de la fibra.
\medskip
\Needspace{10\baselineskip}
\noindent\textbf{rama quark de tipo inferior, profundidad G3.} \(\mathcal B_{d,\mathrm{G3}}\).\par
\noindent\textbf{Dominio de fibra.} Rango 6; celdas 12,14,16,66,68,70; B | G | R; 6 familias de ruta.
\noindent\textbf{Espectro completo de \(K_D\).}
\begin{itemize}
\item \texttt{\detokenize{(40,24,2)x2}}
\item \texttt{\detokenize{(40,78,27)x1}}
\item \texttt{\detokenize{(60,84,28)x1}}
\item \texttt{\detokenize{(80,54,6)x1}}
\item \texttt{\detokenize{(80,66,4)x1}}
\end{itemize}
\noindent\textbf{Espectro completo de \(K_\Omega\).}
\begin{itemize}
\item \texttt{\detokenize{(36,54,4)x1}}
\item \texttt{\detokenize{(80,54,6)x5}}
\end{itemize}
\noindent\textbf{Conclusión en la fibra.} Operador diagonal exacto sobre la fibra \(B_1\); la ruta nominal y el morfismo declarado publican el carácter escalar sin sustituir el espectro de la fibra.
\medskip
\Needspace{10\baselineskip}
\noindent\textbf{rama quark de tipo inferior, profundidad G4.} \(\mathcal B_{d,\mathrm{G4}}\).\par
\noindent\textbf{Dominio de fibra.} Rango 8; celdas 10,18,28,36,46,54,64,72; B | G | R; 8 familias de ruta.
\noindent\textbf{Espectro completo de \(K_D\).}
\begin{itemize}
\item \texttt{\detokenize{(0,24,0)x1}}
\item \texttt{\detokenize{(40,24,2)x1}}
\item \texttt{\detokenize{(40,78,27)x2}}
\item \texttt{\detokenize{(40,84,30)x1}}
\item \texttt{\detokenize{(80,54,7)x1}}
\item \texttt{\detokenize{(80,66,3)x1}}
\item \texttt{\detokenize{(100,66,2)x1}}
\end{itemize}
\noindent\textbf{Espectro completo de \(K_\Omega\).}
\begin{itemize}
\item \texttt{\detokenize{(36,54,4)x1}}
\item \texttt{\detokenize{(80,54,6)x7}}
\end{itemize}
\noindent\textbf{Conclusión en la fibra.} Operador diagonal exacto sobre la fibra \(B_1\); la ruta nominal y el morfismo declarado publican el carácter escalar sin sustituir el espectro de la fibra.
\medskip
\Needspace{10\baselineskip}
\noindent\textbf{rama quark de tipo superior, profundidad G1.} \(\mathcal B_{u,\mathrm{G1}}\).\par
\noindent\textbf{Dominio de fibra.} Rango 4; celdas 31,33,49,51; B | G | R; 3 familias de ruta.
\noindent\textbf{Espectro completo de \(K_D\).}
\begin{itemize}
\item \texttt{\detokenize{(40,54,6)x1}}
\item \texttt{\detokenize{(60,78,25)x1}}
\item \texttt{\detokenize{(80,66,2)x1}}
\item \texttt{\detokenize{(80,66,3)x1}}
\end{itemize}
\noindent\textbf{Espectro completo de \(K_\Omega\).}
\begin{itemize}
\item \texttt{\detokenize{(80,54,6)x4}}
\end{itemize}
\noindent\textbf{Conclusión en la fibra.} Operador diagonal exacto sobre la fibra \(B_1\); la ruta nominal y el morfismo declarado publican el carácter escalar sin sustituir el espectro de la fibra.
\medskip
\Needspace{29\baselineskip}
\noindent\textbf{rama quark de tipo superior, profundidad G3.} \(\mathcal B_{u,\mathrm{G3}}\).\par
\noindent\textbf{Dominio de fibra.} Rango 12; celdas 11,13,15,17,29,35,47,53,65,67,69,71; B | G | R; 12 familias de ruta.
\noindent\textbf{Espectro completo de \(K_D\).}
\begin{itemize}
\item \texttt{\detokenize{(40,24,2)x2}}
\item \texttt{\detokenize{(40,78,27)x2}}
\item \texttt{\detokenize{(60,84,28)x1}}
\item \texttt{\detokenize{(80,54,6)x3}}
\item \texttt{\detokenize{(80,66,3)x1}}
\item \texttt{\detokenize{(80,66,4)x1}}
\item \texttt{\detokenize{(100,66,0)x1}}
\item \texttt{\detokenize{(100,66,2)x1}}
\end{itemize}
\noindent\textbf{Espectro completo de \(K_\Omega\).}
\begin{itemize}
\item \texttt{\detokenize{(56,54,5)x3}}
\item \texttt{\detokenize{(80,54,6)x9}}
\end{itemize}
\noindent\textbf{Conclusión en la fibra.} Operador diagonal exacto sobre la fibra \(B_1\); la ruta nominal y el morfismo declarado publican el carácter escalar sin sustituir el espectro de la fibra.
\medskip
\endgroup


\Needspace{15\baselineskip}
\subsubsection{Morfismos de realización física}

\begingroup
\setlength{\tabcolsep}{4.5pt}
\renewcommand{\arraystretch}{1.28}
\begin{longtable}{@{}>{\raggedright\arraybackslash}p{.14\textwidth}>{\raggedright\arraybackslash}p{.33\textwidth}>{\raggedright\arraybackslash}p{.25\textwidth}>{\raggedright\arraybackslash}p{.22\textwidth}@{}}
\caption{Realización física de los leptones cargados.}\label{tab:masas-realizacion-leptones}\\
\toprule
\textbf{Clase} & \textbf{Objeto HMT} & \textbf{Mapa de realización} & \textbf{Conclusión y condiciones} \\
\midrule
\endfirsthead
\toprule
\textbf{Clase} & \textbf{Objeto HMT} & \textbf{Mapa de realización} & \textbf{Conclusión y condiciones} \\
\midrule
\endhead
\midrule
\multicolumn{4}{r}{\footnotesize Continúa en la página siguiente.}\\
\endfoot
\bottomrule
\endlastfoot
$e^-/e^+$ & Celda central $(5,5)$, ruta $\Gamma_e=(5,5,++)$ y $K_D(\Gamma_e)=K_\Omega(\Gamma_e)=(80,54,6)$. & Reducción de rango uno cerrada. & Clausura beta de rango uno; selector independiente de magnitudes metrológicas empleadas como blanco. \\
$\mu^-/\mu^+$ & familia leptónica cargada; niveles G2; 8 celdas; $K_D$: 7 perfiles; $K_\Omega$: 1 perfil & Generación y firma fina $\to$ ruta persistente $\to\chi_{\rm mass}\to\mathcal L_M$. & Ruta nominal y carácter escalar cerrados; el par operatorial conserva el refinamiento completo de la fibra. \\
$\tau^-/\tau^+$ & familia leptónica cargada; niveles G4; 16 celdas; $K_D$: 8 perfiles; $K_\Omega$: 3 perfiles & Generación y firma fina $\to$ ruta persistente $\to\chi_{\rm mass}\to\mathcal L_M$. & Ruta nominal y carácter escalar cerrados; el par operatorial conserva el refinamiento completo de la fibra. \\
\end{longtable}
\endgroup

\begingroup
\setlength{\tabcolsep}{4.5pt}
\renewcommand{\arraystretch}{1.28}
\begin{longtable}{@{}>{\raggedright\arraybackslash}p{.14\textwidth}>{\raggedright\arraybackslash}p{.33\textwidth}>{\raggedright\arraybackslash}p{.25\textwidth}>{\raggedright\arraybackslash}p{.22\textwidth}@{}}
\caption{Realización física de los neutrinos y transporte de bases.}\label{tab:masas-realizacion-neutrinos}\\
\toprule
\textbf{Clase} & \textbf{Objeto HMT} & \textbf{Mapa de realización} & \textbf{Conclusión y condiciones} \\
\midrule
\endfirsthead
\toprule
\textbf{Clase} & \textbf{Objeto HMT} & \textbf{Mapa de realización} & \textbf{Conclusión y condiciones} \\
\midrule
\endhead
\midrule
\multicolumn{4}{r}{\footnotesize Continúa en la página siguiente.}\\
\endfoot
\bottomrule
\endlastfoot
$\nu_e/\bar\nu_e$ & familia neutra; niveles G1, G2, G3, G4; 20 celdas; $K_D$: 11 perfiles; $K_\Omega$: 5 perfiles & Registro neutro de rango pleno $\to$ autovectores $\to$ base PMNS. & La estructura exterior/CAR fija la fermionicidad; PMNS transporta bases y no crea autovalores de masa. \\
$\nu_\mu/\bar\nu_\mu$ & familia neutra; niveles G1, G2, G3, G4; 20 celdas; $K_D$: 11 perfiles; $K_\Omega$: 5 perfiles & Registro neutro de rango pleno $\to$ autovectores $\to$ base PMNS. & La estructura exterior/CAR fija la fermionicidad; PMNS transporta bases y no crea autovalores de masa. \\
$\nu_\tau/\bar\nu_\tau$ & familia neutra; niveles G1, G2, G3, G4; 20 celdas; $K_D$: 11 perfiles; $K_\Omega$: 5 perfiles & Registro neutro de rango pleno $\to$ autovectores $\to$ base PMNS. & La estructura exterior/CAR fija la fermionicidad; PMNS transporta bases y no crea autovalores de masa. \\
\end{longtable}
\endgroup

\begingroup
\setlength{\tabcolsep}{4.5pt}
\renewcommand{\arraystretch}{1.28}
\begin{longtable}{@{}>{\raggedright\arraybackslash}p{.14\textwidth}>{\raggedright\arraybackslash}p{.33\textwidth}>{\raggedright\arraybackslash}p{.25\textwidth}>{\raggedright\arraybackslash}p{.22\textwidth}@{}}
\caption{Realización física de los quarks: generación, sabor y color.}\label{tab:masas-realizacion-quarks}\\
\toprule
\textbf{Clase} & \textbf{Objeto HMT} & \textbf{Mapa de realización} & \textbf{Conclusión y condiciones} \\
\midrule
\endfirsthead
\toprule
\textbf{Clase} & \textbf{Objeto HMT} & \textbf{Mapa de realización} & \textbf{Conclusión y condiciones} \\
\midrule
\endhead
\midrule
\multicolumn{4}{r}{\footnotesize Continúa en la página siguiente.}\\
\endfoot
\bottomrule
\endlastfoot
$u/\bar u$ & rama quark \emph{up}; niveles G1 y G3; color $\mathbb Z/3$; 16 celdas; $K_D$: 11 perfiles; $K_\Omega$: 2 perfiles & Sabor--color--generación $\to$ ruta persistente $\to\chi_{\rm mass}\to\mathcal L_M$. & Sabor, color y generación determinan la clase de ruta; carácter escalar y par operatorial cerrados. \\
$d/\bar d$ & rama quark \emph{down}; niveles G1, G2, G3, G4; color $\mathbb Z/3$; 20 celdas; $K_D$: 12 perfiles; $K_\Omega$: 4 perfiles & Sabor--color--generación $\to$ ruta persistente $\to\chi_{\rm mass}\to\mathcal L_M$. & Sabor, color y generación determinan la clase de ruta; carácter escalar y par operatorial cerrados. \\
$c/\bar c$ & rama quark \emph{up}; niveles G1 y G3; color $\mathbb Z/3$; 16 celdas; $K_D$: 11 perfiles; $K_\Omega$: 2 perfiles & Sabor--color--generación $\to$ ruta persistente $\to\chi_{\rm mass}\to\mathcal L_M$. & Sabor, color y generación determinan la clase de ruta; carácter escalar y par operatorial cerrados. \\
$s/\bar s$ & rama quark \emph{down}; niveles G1, G2, G3, G4; color $\mathbb Z/3$; 20 celdas; $K_D$: 12 perfiles; $K_\Omega$: 4 perfiles & Sabor--color--generación $\to$ ruta persistente $\to\chi_{\rm mass}\to\mathcal L_M$. & Sabor, color y generación determinan la clase de ruta; carácter escalar y par operatorial cerrados. \\
$t/\bar t$ & rama quark \emph{up}; niveles G1 y G3; color $\mathbb Z/3$; 16 celdas; $K_D$: 11 perfiles; $K_\Omega$: 2 perfiles & Sabor--color--generación $\to$ ruta persistente $\to\chi_{\rm mass}\to\mathcal L_M$. & Sabor, color y generación determinan la clase de ruta; carácter escalar y par operatorial cerrados. \\
$b/\bar b$ & rama quark \emph{down}; niveles G1, G2, G3, G4; color $\mathbb Z/3$; 20 celdas; $K_D$: 12 perfiles; $K_\Omega$: 4 perfiles & Sabor--color--generación $\to$ ruta persistente $\to\chi_{\rm mass}\to\mathcal L_M$. & Sabor, color y generación determinan la clase de ruta; carácter escalar y par operatorial cerrados. \\
\end{longtable}
\endgroup

\Needspace{12\baselineskip}
\paragraph{Localización estructural de los \(6\) sabores de quark}
\begingroup
\setlength{\parskip}{.3\baselineskip}
La posición de un sabor quark se especifica mediante cuatro datos simultáneos: rama, generación, soporte celular y familia de rutas orientadas. Las tres generaciones de una misma rama comparten el soporte del atlas; el clasificador fino de generación conserva la distinción física que todavía debe realizarse como autovector y sección de masa.
\begingroup
\setlength{\tabcolsep}{4.5pt}
\renewcommand{\arraystretch}{1.28}
\begin{longtable}{@{}>{\raggedright\arraybackslash}p{.12\textwidth}>{\raggedright\arraybackslash}p{.17\textwidth}>{\raggedright\arraybackslash}p{.19\textwidth}>{\raggedright\arraybackslash}p{.46\textwidth}@{}}
\caption{Clasificación generacional de los seis sabores quark.}\label{tab:masas-posicion-quarks-seis-sabores}\\
\toprule
\textbf{Sabor} & \textbf{Generación} & \textbf{Rama} & \textbf{Clasificador de rama y generación} \\
\midrule
\endfirsthead
\toprule
\textbf{Sabor} & \textbf{Generación} & \textbf{Rama} & \textbf{Clasificador de rama y generación} \\
\midrule
\endhead
\midrule
\multicolumn{4}{r}{\footnotesize Continúa en la página siguiente.}\\
\endfoot
\bottomrule
\endlastfoot
$u/\bar u$ & primera & superior & $(\operatorname{up},1,\chi_{\rm fina})$; par nominal de ruta $(21,30)$ sobre el eje $\mathrm{N\!-\!S}$ \\
$d/\bar d$ & primera & inferior & $(\operatorname{down},1,\chi_{\rm fina})$; par nominal de ruta $(20,32)$ sobre el eje $\mathrm{E\!-\!O}$ \\
$c/\bar c$ & segunda & superior & $(\operatorname{up},2,\chi_{\rm fina})$; par nominal de ruta $(21,29)$ sobre el eje $\mathrm{E\!-\!O}$ \\
$s/\bar s$ & segunda & inferior & $(\operatorname{down},2,\chi_{\rm fina})$; par nominal de ruta $(17,27)$ sobre el eje $\mathrm{N\!-\!S}$ \\
$t/\bar t$ & tercera & superior & $(\operatorname{up},3,\chi_{\rm fina})$; par nominal de ruta $(21,26)$ sobre el eje $\mathrm{N\!-\!S}$ \\
$b/\bar b$ & tercera & inferior & $(\operatorname{down},3,\chi_{\rm fina})$; par nominal de ruta $(21,30)$ sobre el eje $\mathrm{E\!-\!O}$ \\
\end{longtable}
\endgroup

\Needspace{18\baselineskip}
\noindent\textbf{Soporte estructural de la rama superior.}\par
\noindent\textit{Multisecciones.} \(\mathcal B_{u,\mathrm{G1}}\), \(\mathcal B_{u,\mathrm{G3}}\).\par
\noindent\textit{Profundidades y cardinales.} $\mathrm{G1}$, $\mathrm{G3}$; 16 celdas y 14 familias de ruta distintas.\par
\noindent\textit{Órbita cromática y orientaciones.} $\mathcal C=\{B,G,R\}$ y $\mathcal O=\{(-1,-1),(-1,1),(1,-1),(1,1)\}$.\par
\noindent\textit{Celdas y familias orientadas $B_1$.} Cada cadena $k\leftrightarrow(r,c)\leftrightarrow\Gamma_{r,c}^{(h,v)}$ conserva el índice celular, su coordenada en la carta $9\times9$ y los dos signos de orientación del registro de ruta enriquecido.\par
\noindent Fila $r=2$: $11\leftrightarrow(2,2)\leftrightarrow\Gamma_{2,2}^{(-,+)}$, $13\leftrightarrow(2,4)\leftrightarrow\Gamma_{2,4}^{(+,-)}$, $15\leftrightarrow(2,6)\leftrightarrow\Gamma_{2,6}^{(-,+)}$, $17\leftrightarrow(2,8)\leftrightarrow\Gamma_{2,8}^{(-,-)}$.\par
\noindent Fila $r=4$: $29\leftrightarrow(4,2)\leftrightarrow\Gamma_{4,2}^{(+,-)}$, $31\leftrightarrow(4,4)\leftrightarrow\Gamma_{4,4}^{(-,+)}$, $33\leftrightarrow(4,6)\leftrightarrow\Gamma_{4,6}^{(+,+)}$, $35\leftrightarrow(4,8)\leftrightarrow\Gamma_{4,8}^{(-,-)}$.\par
\noindent Fila $r=6$: $47\leftrightarrow(6,2)\leftrightarrow\Gamma_{6,2}^{(-,+)}$, $49\leftrightarrow(6,4)\leftrightarrow\Gamma_{6,4}^{(-,-)}$, $51\leftrightarrow(6,6)\leftrightarrow\Gamma_{6,6}^{(+,+)}$, $53\leftrightarrow(6,8)\leftrightarrow\Gamma_{6,8}^{(+,+)}$.\par
\noindent Fila $r=8$: $65\leftrightarrow(8,2)\leftrightarrow\Gamma_{8,2}^{(+,+)}$, $67\leftrightarrow(8,4)\leftrightarrow\Gamma_{8,4}^{(+,+)}$, $69\leftrightarrow(8,6)\leftrightarrow\Gamma_{8,6}^{(-,-)}$, $71\leftrightarrow(8,8)\leftrightarrow\Gamma_{8,8}^{(+,+)}$.\par
\noindent\textit{Morfismo de realización cerrado.} Para $g\in\{1,2,3\}$, la rama, la generación, la órbita de color y la firma fina determinan
\[\chi_{\rm fina}(\operatorname{up},g,\mathcal C)\longrightarrow\Gamma_{\rm enr}\longrightarrow\bigl(\operatorname{quot}_{\mathcal C},S(B_D,B_\Omega)\bigr).\]
Esta flecha selecciona el autovector físico y publica después la coordenada escalar en $\mathcal L_M$ sin usar el valor de contraste.\par
\medskip
\Needspace{18\baselineskip}
\noindent\textbf{Soporte estructural de la rama inferior.}\par
\noindent\textit{Multisecciones.} \(\mathcal B_{d,\mathrm{G1}}\), \(\mathcal B_{d,\mathrm{G2}}\), \(\mathcal B_{d,\mathrm{G3}}\), \(\mathcal B_{d,\mathrm{G4}}\).\par
\noindent\textit{Profundidades y cardinales.} $\mathrm{G1}$, $\mathrm{G2}$, $\mathrm{G3}$, $\mathrm{G4}$; 20 celdas y 18 familias de ruta distintas.\par
\noindent\textit{Órbita cromática y orientaciones.} $\mathcal C=\{B,G,R\}$ y $\mathcal O=\{(-1,-1),(-1,1),(1,-1),(1,1)\}$.\par
\noindent\textit{Celdas y familias orientadas $B_1$.} Cada cadena $k\leftrightarrow(r,c)\leftrightarrow\Gamma_{r,c}^{(h,v)}$ conserva el índice celular, su coordenada en la carta $9\times9$ y los dos signos de orientación del registro de ruta enriquecido.\par
\noindent Fila $r=2$: $10\leftrightarrow(2,1)\leftrightarrow\Gamma_{2,1}^{(+,+)}$, $12\leftrightarrow(2,3)\leftrightarrow\Gamma_{2,3}^{(-,-)}$, $14\leftrightarrow(2,5)\leftrightarrow\Gamma_{2,5}^{(+,+)}$, $16\leftrightarrow(2,7)\leftrightarrow\Gamma_{2,7}^{(-,+)}$, $18\leftrightarrow(2,9)\leftrightarrow\Gamma_{2,9}^{(-,+)}$.\par
\noindent Fila $r=4$: $28\leftrightarrow(4,1)\leftrightarrow\Gamma_{4,1}^{(-,-)}$, $30\leftrightarrow(4,3)\leftrightarrow\Gamma_{4,3}^{(+,+)}$, $32\leftrightarrow(4,5)\leftrightarrow\Gamma_{4,5}^{(+,-)}$, $34\leftrightarrow(4,7)\leftrightarrow\Gamma_{4,7}^{(-,+)}$, $36\leftrightarrow(4,9)\leftrightarrow\Gamma_{4,9}^{(-,+)}$.\par
\noindent Fila $r=6$: $46\leftrightarrow(6,1)\leftrightarrow\Gamma_{6,1}^{(-,-)}$, $48\leftrightarrow(6,3)\leftrightarrow\Gamma_{6,3}^{(-,+)}$, $50\leftrightarrow(6,5)\leftrightarrow\Gamma_{6,5}^{(-,+)}$, $52\leftrightarrow(6,7)\leftrightarrow\Gamma_{6,7}^{(-,+)}$, $54\leftrightarrow(6,9)\leftrightarrow\Gamma_{6,9}^{(+,-)}$.\par
\noindent Fila $r=8$: $64\leftrightarrow(8,1)\leftrightarrow\Gamma_{8,1}^{(+,-)}$, $66\leftrightarrow(8,3)\leftrightarrow\Gamma_{8,3}^{(-,+)}$, $68\leftrightarrow(8,5)\leftrightarrow\Gamma_{8,5}^{(-,+)}$, $70\leftrightarrow(8,7)\leftrightarrow\Gamma_{8,7}^{(+,-)}$, $72\leftrightarrow(8,9)\leftrightarrow\Gamma_{8,9}^{(-,+)}$.\par
\noindent\textit{Morfismo de realización cerrado.} Para $g\in\{1,2,3\}$, la rama, la generación, la órbita de color y la firma fina determinan
\[\chi_{\rm fina}(\operatorname{down},g,\mathcal C)\longrightarrow\Gamma_{\rm enr}\longrightarrow\bigl(\operatorname{quot}_{\mathcal C},S(B_D,B_\Omega)\bigr).\]
Esta flecha selecciona el autovector físico y publica después la coordenada escalar en $\mathcal L_M$ sin usar el valor de contraste.\par
\medskip
\endgroup

\begingroup
\setlength{\tabcolsep}{4.5pt}
\renewcommand{\arraystretch}{1.28}
\begin{longtable}{@{}>{\raggedright\arraybackslash}p{.14\textwidth}>{\raggedright\arraybackslash}p{.33\textwidth}>{\raggedright\arraybackslash}p{.25\textwidth}>{\raggedright\arraybackslash}p{.22\textwidth}@{}}
\caption{Realización física de los sectores gauge y escalar.}\label{tab:masas-realizacion-gauge-escalar}\\
\toprule
\textbf{Clase} & \textbf{Objeto HMT} & \textbf{Mapa de realización} & \textbf{Conclusión y condiciones} \\
\midrule
\endfirsthead
\toprule
\textbf{Clase} & \textbf{Objeto HMT} & \textbf{Mapa de realización} & \textbf{Conclusión y condiciones} \\
\midrule
\endhead
\midrule
\multicolumn{4}{r}{\footnotesize Continúa en la página siguiente.}\\
\endfoot
\bottomrule
\endlastfoot
fotón $\gamma$ & Carta nula sectorial y representación gauge correspondiente. & Inclusión exacta de la sección nula en la recta de masa. & Nulidad exacta en su carta sectorial. \\
gluones $g$ & Carta nula sectorial y representación gauge correspondiente. & Inclusión exacta de la sección nula en la recta de masa. & Nulidad exacta en su carta sectorial. \\
$W^\pm$ & ruta CT--108 W: 15/37, eje E-O, $(N,E,S,O)=(23,44,11,30)$, hoja/residuo $(41,8)$, $(n_A,s_C)=(94,-1)$, $(W_+,W_-)=(563,565)$; carácter escalar exacto & Ruta CT--108 nominal $\to(n_A,s_C)\to(W_+,W_-)\to\chi_{\rm mass}\to\mathcal L_M$. & Ruta gauge CT--108, firma y carácter escalar determinados antes del contraste. \\
$Z^0$ & ruta CT--108 Z: 20/23, eje E-O, $(N,E,S,O)=(33,45,10,20)$, hoja/residuo $(57,7)$, $(n_A,s_C)=(95,-1)$, $(W_+,W_-)=(569,571)$; carácter escalar exacto & Ruta CT--108 nominal $\to(n_A,s_C)\to(W_+,W_-)\to\chi_{\rm mass}\to\mathcal L_M$. & Ruta gauge CT--108, firma y carácter escalar determinados antes del contraste. \\
Higgs $H$ & ruta CT--108 H: 24/29, eje N-S, $(N,E,S,O)=(40,34,22,12)$, hoja/residuo $(47,10)$, $(n_A,s_C)=(97,9)$, $(W_+,W_-)=(591,573)$; carácter escalar exacto & Ruta CT--108 nominal $\to(n_A,s_C)\to(W_+,W_-)\to\chi_{\rm mass}\to\mathcal L_M$. & Ruta escalar CT--108, firma y carácter determinados antes del contraste. \\
\end{longtable}
\endgroup

\begingroup
\setlength{\tabcolsep}{4.5pt}
\renewcommand{\arraystretch}{1.28}
\begin{longtable}{@{}>{\raggedright\arraybackslash}p{.14\textwidth}>{\raggedright\arraybackslash}p{.33\textwidth}>{\raggedright\arraybackslash}p{.25\textwidth}>{\raggedright\arraybackslash}p{.22\textwidth}@{}}
\caption{Realización de compuestos, sectores topológicos y estados virtuales.}\label{tab:masas-realizacion-no-escalar}\\
\toprule
\textbf{Clase} & \textbf{Objeto HMT} & \textbf{Mapa de realización} & \textbf{Conclusión y condiciones} \\
\midrule
\endfirsthead
\toprule
\textbf{Clase} & \textbf{Objeto HMT} & \textbf{Mapa de realización} & \textbf{Conclusión y condiciones} \\
\midrule
\endhead
\midrule
\multicolumn{4}{r}{\footnotesize Continúa en la página siguiente.}\\
\endfoot
\bottomrule
\endlastfoot
mesones & Registro doble, frontera de enlace y firma compuesta. & Registro compuesto y ligadura $\to$ carta escalar. & Composición de registros; la ligadura precede a la carta escalar. \\
bariones & Registro triple, acarreo y firma compuesta. & Registro compuesto y ligadura $\to$ carta escalar. & Composición de registros; la ligadura precede a la carta escalar. \\
Fibonacci $(1,\tau)$ & Anillo de fusión Fibonacci: $\tau\otimes\tau=1+\tau$. & Fusión y trenzado $\to$ observable físico del sector. & Codominio exacto de fusión y trenzado; el contraste usa sus invariantes propios. \\
Ising $(1,\psi,\sigma)$; realización Majorana separada & Anillo de fusión de Ising: $(1,\psi,\sigma)$. & Fusión y trenzado $\to$ observable físico del sector. & Codominio exacto de fusión y trenzado; el contraste usa sus invariantes propios. \\
sectores $\mathbb Z/6$ & Carácter de trenzado sobre $\mathbb Z/6$. & Fusión y trenzado $\to$ observable físico del sector. & Codominio exacto de fusión y trenzado; el contraste usa sus invariantes propios. \\
secciones virtuales & Sistema inverso de extensiones y horizonte de persistencia. & Prolongación registrada $\to$ horizonte físico de persistencia. & Objeto exacto de prolongación y obstrucción; su salida es el horizonte de persistencia. \\
\end{longtable}
\endgroup


\Needspace{15\baselineskip}
\subsubsection{Comparabilidad física en los codominios nativos}

La comparación externa se define sobre el codominio nativo de cada clase. En las
clases con sección escalar cerrada se contrastan magnitudes homogéneas, con unidad,
esquema y escala declarados. En las fibras operatoriales se contrasta el espectro
interno y se conserva de forma explícita el morfismo de realización cerrado que
selecciona la ruta sin usar el valor externo. Los registros mesónicos y bariónicos sólo admiten comparación estado por
estado después de incorporar la frontera de ligadura. En el sector Fibonacci están
cerrados el anillo basado y \(\operatorname{FPdim}(\tau)=\varphi\); las matrices
\(F/R\), el pentágono, el hexágono, el Hamiltoniano y la brecha pertenecen a la
realización posterior. Ising y \(\mathbb Z/6\) conservan por separado sus invariantes
de fusión y trenzado, y la virtualidad se compara mediante persistencia y estructura
analítica. En ninguno de estos seis
casos existe una masa escalar universal de clase; asignarla sustituiría el objeto
demostrado por una cifra sin morfismo.
